% =============================================================================
% Section 9 -- Conclusion
% =============================================================================
\section{Conclusion}
\label{sec:conclusion}

We have introduced the Agent Performance Improvement Plan (APIP) as a
structured post-deployment remediation framework for customer-facing AI agents.
The framework adapts the human PIP's lifecycle primitives---trigger, cause
attribution, tiered intervention, remediation window, exit evaluation---while
acknowledging that agents cannot self-direct improvement. We ground the
evaluation discipline in Kirkpatrick's four-level model to distinguish durable
recovery from superficial metric improvement.

The APIP framework addresses a genuine gap in current practice: the absence of
structured, documented, and measurable remediation protocols for
underperforming production agents. Through a hypothetical case study, we have
shown that unstructured remediation creates avoidable regression risk through
failure mode misclassification, and that the APIP's cause attribution phase
provides a direct mitigation. The formal schema provides a foundation for
tooling integration and the audit trail documentation increasingly required by
AI governance frameworks~\cite{brundage2020trustworthy,eu2024aiact,wachter2017counterfactual}.

We have identified behavioral contract specification, scalable causal
attribution, multi-agent mesh extension, and organizational ownership as the
primary open problems for future work. The mesh disruption problem in
particular---where a locally high-performing new agent degrades peer agent
performance through shared context space dynamics---reveals a causal
attribution failure mode that motivates a companion extension of this framework
to multi-agent deployments. By grounding mesh disruption analysis in Event
System Theory~\cite{morgeson2015event}, team membership change
research~\cite{summers2012team,trainer2020membership}, Transactive Memory
Systems~\cite{ren2011transactive,wegner1987transactive}, and empirical
multi-agent scaling findings~\cite{kim2025scaling}, we bridge organizational
behavior and multi-agent AI. The APIP serves as the remediation complement to
drift detection (Rath~\cite{rath2026drift}) and production diagnosis (Pan et
al.~\cite{pan2025measuring}), and as the protocol layer above run-level repair
tooling~\cite{zhu2026agentdebugx}: detection identifies that an agent is
failing; the APIP governs what to do next.

Our central argument is simple: deployed agents should be treated as subjects
of ongoing accountability rather than as static shipped artifacts. The
discipline of structured remediation---trigger, diagnose, intervene,
evaluate---is not new. It is well established in software reliability
engineering, in organizational management, and in clinical practice. The AI
field has the opportunity to adopt it deliberately, before regulatory
requirements make adoption mandatory.
